% -*- coding: utf-8 -*-
% !TEX program = xelatex
\documentclass[plain,noanswer,medmath]{../../template/jnuexam}
\usepackage{../../template/kysx}

\begin{document}


\examtitle{2002}{2}


\loadproblems{1}{2002P1}%载入试卷一


\makepart{填空题}{本题共5小题，每小题3分，满分15分}


\begin{problem}
设函数$f(x) = \begin{cases}
  \frac{1 - \e^{\tan x}}{\arcsin \frac{x}{2}},& x > 0 \\
  a\e^{2x}, & x \le 0
\end{cases}$在$x = 0$处连续，则$a =$ \fillin{}．
\end{problem}


\begin{problem}
位于曲线$y = x\e^{- x}(0 \le x < +\infty)$下方，$x$轴上方的无界图形的面积是 \fillin{}．
\end{problem}


\useproblem{1}{1}{3}%【同试卷一第一(3)题】


\begin{problem}
$\lim_{n \to \infty} \frac{1}{n}\left[\sqrt{1 + \cos \frac{\pi}{n} +} \sqrt{1 + \cos \frac{2\pi}{n} +}
  \cdots + \sqrt{1 + \cos \frac{n\pi}{n}} \right] =$ \fillin{}．
\end{problem}


\begin{problem}
矩阵$\begin{pmatrix}
   0 & -2 & -2 \\
   2 &  2 & -2 \\
  -2 & -2 & 2
\end{pmatrix}$的非零特征值是 \fillin{}．
\end{problem}


\makepart{选择题}{本题共5小题，每小题3分，共15分}


\begin{problem}
设函数$f(u)$可导，$y = f(x^2)$当自变量$x$在$x = - 1$处取得增量$\Delta x = - 0.1$时，
相应的函数增量$\Delta y$的线性主部为$0.1$，则$f'(1) =$\pickout{}
\begin{abcd}
\item $-1$．
\item $0.1$．
\item $1$．
\item $0.5$．
\end{abcd}
\end{problem}


\begin{problem}
设函数$f(x)$连续，则下列函数中，必为偶函数的是\pickout{}
\begin{abcd}
\item $\int_0^x f(t^2) \dt$．
\item $\int_0^x f^2(t) \dt$．
\item $\int_0^x t[f(t) - f(- t)]\dt$．
\item $\int_0^x t[f(t) + f(- t)]\dt$．
\end{abcd}
\end{problem}


\begin{problem}
设$y = (x)$是二阶常系数微分方程$y'' + py' + qy = \e^{3x}$满足初始条$y(0) = y'(0) = 0$的特解，
则当$x \to 0$，函数$\frac{\ln(1 + x^2)}{y(x)}$的极限\pickout{}
\begin{abcd}
\item 不存在．
\item 等于$1$．
\item 等于$2$．
\item 等于$3$．
\end{abcd}
\end{problem}


\useproblem{1}{2}{3}%【同试卷一第二(3)题】


\begin{problem}
设向量组$\alpha_1,\alpha_2,\alpha_3$线性无关，向量$\beta_1$可由$\alpha_1,\alpha_2,\alpha_3$线性表示，
而向量$\beta_2$不能由$\alpha_1,\alpha_2,\alpha_3$线性表示，则对于任意常数$k$，必有\pickout{}
\begin{abcd}
\item $\alpha_1,\alpha_2,\alpha_3$,$k\beta_1 + \beta_2$线性无关．
\item $\alpha_1,\alpha_2,\alpha_3$,$k\beta_1 + \beta_2$线性相关．
\item $\alpha_1,\alpha_2,\alpha_3$,$\beta_1 + k\beta_2$线性无关．
\item $\alpha_1,\alpha_2,\alpha_3$,$\beta_1 + k\beta_2$线性相关．
\end{abcd}
\end{problem}


\makepart{}{本题满分6分}

\begin{problem}
已知曲线的极坐标方程是$r = 1 - \cos\theta$，
求该曲线上对应于$\theta = \frac{\pi}{6}$处的切线与法线的直角坐标方程．
\end{problem}


\makepart{}{本题满分7分}

\begin{problem}
设$f(x) = \begin{cases}
         2x + \frac{3}{2}x^2, & -1 \le x < 0; \\
  \frac{x\e^x}{(\e^x + 1)^2}, & 0 \le x \le 1.
\end{cases}$求函数$F(x) = \int_{-1}^x f(t)\dt$的表达式．
\end{problem}


\makepart{}{本题满分7分}

\begin{problem}
已知函数$f(x)$在$(0, +\infty)$内可导$f(x) > 0$，$\lim_{x \to +\infty} f(x) = 1$，且满足
$$\lim_{h \to 0} \left(\frac{f(x + hx)}{f(x)} \right)^{\frac{1}{h}} = \e^{\frac{1}{x}},$$
求$f(x)$．
\end{problem}


\makepart{}{本题满分8分}

\begin{problem}
求微分方程$x\dy + (x - 2y)\dx = 0$的一个解$y = y(x)$，
使得由曲线$y = y(x)$与直线$x = 1$，$x = 2$以及$x$轴所围成的平面图形绕$x$轴旋转一周的旋转体体积最小．
\end{problem}


\makepart{}{本题满分7分}

\begin{problem}
某闸门的形状与大小如图所示，其中直线$l$为对称轴，闸门的上部为矩形$ABCD$，下部由二次抛物线与线段$AB$所围成，
当水面与闸门的上端相平时，欲使闸门矩形部分承受的水压力与闸门下部承受的水压力之比为$5:4$，
闸门矩形部分的高$h$应为多少$m$(米)？
\par\centerline{\includegraphics[width=0.4\linewidth]{2002B1}}
\end{problem}


\makepart{}{本题满分8分}

\begin{problem}
设$0 < x_1 < 3$，$x_{n + 1} = \sqrt{x_n(3 - x_n)}$（$n = 1,2, \cdots$），
证明数列$\left\{x_n \right\}$的极限存在，并求此极限．
\end{problem}


\makepart{}{本题满分8分}

\begin{problem}
设$0 < a < b$，证明不等式$\frac{2a}{a^2 + b^2} < \frac{\ln b - \ln a}{b - a} < \frac{1}{\sqrt{ab}}$．
\end{problem}


\makepart{}{本题满分8分}

\begin{problem}
设函数$f(x)$在$x = 0$的某邻域内具有二阶连续导数，且$f(0) \ne 0$，$f'(0) \ne 0$，$f''(0) \ne 0$．
证明：存在惟一的一组实数$\lambda_1,\lambda_2,\lambda_3$，使得当$h \to 0$时，
$$\lambda_1 f(h) + \lambda_2 f(2h) + \lambda_3 f(3h) - f(0)$$
是比$h^2$高阶的无穷小．
\end{problem}


\makepart{}{本题满分6分}

\begin{problem}
已知$A,B$为$3$阶矩阵，且满足$2A^{-1} B = B - 4E$，其中$E$是$3$阶单位矩阵．
\begin{enumerate*}
\item 证明：矩阵$A - 2E$可逆；
\item 若$B = \begin{pmatrix}
  1 & -2 & 0 \\
  1 &  2 & 0 \\
  0 &  0 & 2
\end{pmatrix}$，求矩阵$A$．
\end{enumerate*}
\end{problem}


\makepart{}{本题满分6分}%【同试卷一第九题】

\useproblem{1}{9}{0}


\end{document}
